\section{Introduction}

A robot's body and its control objective are fundamentally coupled. Lengthen a leg and the optimal gait changes. Tweak the reward to penalize energy and an otherwise agile body crawls. This coupling implies that optimizing morphology under a fixed reward can yield bodies that become fragile when the reward changes, while engineering rewards for a fixed body limits achievable performance. The joint search space is high-dimensional and nonlinear, making exhaustive exploration impractical. Yet addressing this challenge is critical as co-designed robots can outperform their fixed-body or fixed-reward counterparts by exploiting body--controller synergies that neither optimization can discover~\cite{sims1994evolving}.

Recent work has shown that large language models (LLMs) can generate reward code~\cite{ma2024eureka, ma2024dreureka} or propose morphology edits~\cite{qiu2026robomorph, ringel2025text2robot}, but most existing pipelines treat body and reward as separate problems. This misses opportunities for co-adaptation, where a reward tuned for one morphology may be suboptimal for another. Worse, single-agent generation tends to reuse familiar patterns, narrowing exploration precisely where co-design demands diversity.

We introduce \method{} (\abbrev{}), a multi-agent LLM framework that addresses these limitations through structured debate (Figure~\ref{fig:d2c}). A \emph{design agent} proposes morphology edits while a \emph{control agent} critiques designs and writes reward code, engaging in a thesis--antithesis--synthesis loop where proposals are refined before evaluation. A panel of criterion-specific LLM judges, each focused on a different objective, provides feedback that steers exploration away from narrow optima. All selection decisions are grounded in physics simulation, where candidates are trained in Brax~\cite{brax2021github} and ranked by a standardized task score, ensuring that debate rhetoric is governed by empirical performance.

On five MuJoCo locomotion tasks, \abbrev{} achieves $3.2\times$ the score of the unmodified Ant morphology and the highest default-normalized score among the evaluated LLM-based and black-box baselines under the same evaluation metric. Cross-over experiments reveal that \abbrev{}-generated rewards transfer to default morphologies, improving 4 of 5 tasks without design changes. This suggests the learned shaping captures transferable locomotion principles rather than morphology-specific hacks~\cite{ng1999policy}.

Our contributions are as follows:
\begin{enumerate}[leftmargin=*,nosep]
    \item \textbf{A multi-agent debate framework for co-design.} We introduce \abbrev{}, which, to our knowledge, is the first framework that jointly optimizes morphology and reward through a role-separated thesis--antithesis--synthesis debate loop with cross-agent critique.
    \item \textbf{Criterion-specific LLM judges.} Role-specialized judges provide multi-objective critiques that improve metric coverage and steer exploration toward diverse solutions.
    \item \textbf{Ablation and transfer experiments.} We evaluate \abbrev{} against LLM-based and black-box baselines on MuJoCo locomotion tasks, with cross-over experiments showing reward transferability to unmodified morphologies and ablations quantifying the 18--35\% gains from iterative debate over single-pass generation.
\end{enumerate}
